\documentclass[11pt,a4paper]{article}
\newcommand{\gradelabel}{Grade 9}
\newcommand{\leconhdr}{Feu tricolore automatisé}
% Préambule commun aux plans de cours de robotique (Grades 7 à 12).
% Chaque plan de cours commence par \documentclass[11pt,a4paper]{article}, définit
% \gradelabel et \leconhdr (en-têtes de page), puis % Préambule commun aux plans de cours de robotique (Grades 7 à 12).
% Chaque plan de cours commence par \documentclass[11pt,a4paper]{article}, définit
% \gradelabel et \leconhdr (en-têtes de page), puis % Préambule commun aux plans de cours de robotique (Grades 7 à 12).
% Chaque plan de cours commence par \documentclass[11pt,a4paper]{article}, définit
% \gradelabel et \leconhdr (en-têtes de page), puis \input{../../preambule_robotique.tex}.
\usepackage[T1]{fontenc}
\usepackage[utf8]{inputenc}
\usepackage{lmodern}
\usepackage[french]{babel}
\usepackage[margin=2cm,headheight=14pt]{geometry}
\usepackage[most]{tcolorbox}
\usepackage{xcolor,tabularx,booktabs,array,enumitem,fancyhdr,amssymb,tikz}
\usetikzlibrary{arrows.meta,positioning}

\pagestyle{fancy}
\fancyhf{}
\lhead{\small Robotique -- \gradelabel{} -- Novembre 2026}
\rhead{\small \leconhdr}
\cfoot{\thepage}
\setlist{itemsep=1pt,topsep=2pt,leftmargin=*}
\renewcommand{\arraystretch}{1.35}
\newcolumntype{Y}{>{\raggedright\arraybackslash}X}
\newcommand{\cb}{$\square$\,}
\newcommand{\fl}{$\rightarrow$}
\newcommand{\lig}{\noindent\rule{\linewidth}{0.3pt}\\[5mm]}

\newtcolorbox{etape}[3]{enhanced,breakable,colback=#1!6,colframe=#1!70!black,
  coltitle=white,fonttitle=\bfseries,title={#2\hfill #3},boxrule=0.8pt,arc=2mm,
  before skip=8pt,after skip=8pt}
\newtcolorbox{astuce}[1][Différenciation et besoins particuliers]{breakable,
  colback=yellow!12,colframe=yellow!50!black,fonttitle=\bfseries\small,
  title={#1},boxrule=0.5pt,arc=1mm,before skip=6pt,after skip=6pt}
\newtcolorbox{code}[1][Code]{breakable,colback=black!4,colframe=black!60,
  fonttitle=\bfseries\small,title={#1},boxrule=0.5pt,arc=1mm,fontupper=\ttfamily\small,
  before skip=6pt,after skip=6pt}
% Blocs de code coloré (écrire le code sans accents : comme dans l'éditeur)
\newtcblisting{arduino}[1][Code Arduino]{listing only,breakable,colback=black!4,
  colframe=black!60,fonttitle=\bfseries\small,title={#1},boxrule=0.5pt,arc=1mm,
  before skip=6pt,after skip=6pt,
  listing options={language=C++,basicstyle=\ttfamily\small,columns=fullflexible,
  keepspaces=true,showstringspaces=false,commentstyle=\color{green!40!black}\itshape,
  keywordstyle=\color{blue!70!black}\bfseries}}
\newtcblisting{python}[1][Code Python]{listing only,breakable,colback=black!4,
  colframe=black!60,fonttitle=\bfseries\small,title={#1},boxrule=0.5pt,arc=1mm,
  before skip=6pt,after skip=6pt,
  listing options={language=Python,basicstyle=\ttfamily\small,columns=fullflexible,
  keepspaces=true,showstringspaces=false,commentstyle=\color{green!40!black}\itshape,
  keywordstyle=\color{blue!70!black}\bfseries}}
\newcommand{\entete}[2]{%
  \begin{tcolorbox}[colback=#1!10,colframe=#1!70!black,boxrule=0.8pt,arc=2mm]
  {\Large\bfseries #2}\hfill Nom : \rule{4cm}{0.3pt}\quad Date : \rule{2.2cm}{0.3pt}
  \end{tcolorbox}}

% Bandeau de titre : #1 titre, #2 semaine
\newcommand{\titreblock}[2]{%
\begin{tcolorbox}[colback=teal!10,colframe=teal!70!black,arc=2mm,boxrule=1pt]
\centering
{\LARGE\bfseries #1}\\[2mm]
{\large Robotique -- \gradelabel{} -- Plan de cours de 2 heures}\\[1mm]
Enseignant : Francis Ki \quad|\quad #2
\end{tcolorbox}}

% Tableau du minutage : 6 contenus (accroche, mini-leçon, guidée, autonome, mise en commun, exit ticket)
\newcommand{\minutage}[6]{%
\begin{tabularx}{\linewidth}{>{\bfseries}p{2.2cm}p{2.6cm}Y p{3cm}}
\toprule
Minutes & Étape & Contenu & Organisation\\
\midrule
0 -- 10 & Accroche & #1 & Classe entière\\
10 -- 30 & Mini-leçon & #2 & Classe entière\\
30 -- 60 & Pratique guidée & #3 & Collectif puis binômes\\
60 -- 95 & Pratique autonome & #4 & Individuel sur poste\\
95 -- 110 & Mise en commun & #5 & Classe entière\\
110 -- 120 & Exit ticket & #6 & Individuel sur poste\\
\bottomrule
\end{tabularx}}

% Étayage progressif : #1 aide 1, #2 aide 2, #3 aide 3
\newcommand{\aides}[3]{%
\textbf{Étayage progressif de l'aide.}
\begin{center}
\begin{tabularx}{\linewidth}{>{\bfseries}p{2.4cm}Y p{4.4cm}}
\toprule
Niveau d'aide & Contenu & Qui intervient\\
\midrule
Aide 1 & #1 & Le binôme seul\\
Aide 2 & #2 & Le binôme avec la carte d'aide\\
Aide 3 & #3 & L'enseignant\\
\bottomrule
\end{tabularx}
\end{center}}

% Les trois groupes de la pratique autonome : #1 tâche A, #2 tâche B, #3 tâche C
\newcommand{\groupes}[3]{%
\textbf{Consigne commune\,:} chaque élève ouvre le fichier de son groupe, le renomme avec son prénom et travaille seul sur son poste. L'enseignant passe d'abord auprès du groupe A pendant 10 min.

\medskip
\begin{tabularx}{\linewidth}{>{\bfseries}p{3cm}Y p{2.2cm}}
\toprule
Groupe & Tâche sur poste & Annexe\\
\midrule
A : en dessous du niveau & #1 & Annexe C\\
B : au niveau & #2 & Annexe D\\
C : au-dessus du niveau & #3 & Annexe E\\
\bottomrule
\end{tabularx}}

% Mise en commun : #1 groupe A, #2 groupe B, #3 groupe C, #4 conclusion
\newcommand{\miseencommun}[4]{%
\begin{etape}{teal}{5. Mise en commun}{95 -- 110 min}
\begin{enumerate}
 \item (2 min) Retour au calme, écrans mis en veille.
 \item (9 min) Trois passages courts de 3 min, avec valorisation :
 \begin{itemize}
  \item \textbf{Groupe A} : #1
  \item \textbf{Groupe B} : #2
  \item \textbf{Groupe C} : #3
 \end{itemize}
 \item (4 min) Conclusion : #4
\end{enumerate}
\end{etape}}

% Clôture : #1 question 1, #2 question 2
\newcommand{\cloture}[2]{%
\begin{etape}{gray}{6. Clôture : exit ticket}{110 -- 120 min}
Formulaire en ligne ou papier (annexe G) :
\begin{enumerate}
 \item #1
 \item #2
 \item Auto-évaluation : \emph{Je maîtrise / Je progresse / J'ai besoin d'aide.}
\end{enumerate}
Lire les réponses avant la séance suivante pour ajuster les groupes de besoin.
\end{etape}}

% Grille d'observation : #1 #2 #3 intitulés de colonnes
\newcommand{\grille}[3]{%
\subsection*{Grille d'observation (enseignant)}
\begin{tabularx}{\linewidth}{Y p{2.4cm} p{2.4cm} p{2.4cm} Y}
\toprule
Élève & #1 & #2 & #3 & Groupe suivant\\
\midrule
 & & & & \\ & & & & \\ & & & & \\ & & & & \\ & & & & \\
\bottomrule
\end{tabularx}}

% Signature de l'exit ticket
\newcommand{\autoeval}{Je me sens : \cb Je maîtrise \quad \cb Je progresse \quad \cb J'ai besoin d'aide}
.
\usepackage[T1]{fontenc}
\usepackage[utf8]{inputenc}
\usepackage{lmodern}
\usepackage[french]{babel}
\usepackage[margin=2cm,headheight=14pt]{geometry}
\usepackage[most]{tcolorbox}
\usepackage{xcolor,tabularx,booktabs,array,enumitem,fancyhdr,amssymb,tikz}
\usetikzlibrary{arrows.meta,positioning}

\pagestyle{fancy}
\fancyhf{}
\lhead{\small Robotique -- \gradelabel{} -- Novembre 2026}
\rhead{\small \leconhdr}
\cfoot{\thepage}
\setlist{itemsep=1pt,topsep=2pt,leftmargin=*}
\renewcommand{\arraystretch}{1.35}
\newcolumntype{Y}{>{\raggedright\arraybackslash}X}
\newcommand{\cb}{$\square$\,}
\newcommand{\fl}{$\rightarrow$}
\newcommand{\lig}{\noindent\rule{\linewidth}{0.3pt}\\[5mm]}

\newtcolorbox{etape}[3]{enhanced,breakable,colback=#1!6,colframe=#1!70!black,
  coltitle=white,fonttitle=\bfseries,title={#2\hfill #3},boxrule=0.8pt,arc=2mm,
  before skip=8pt,after skip=8pt}
\newtcolorbox{astuce}[1][Différenciation et besoins particuliers]{breakable,
  colback=yellow!12,colframe=yellow!50!black,fonttitle=\bfseries\small,
  title={#1},boxrule=0.5pt,arc=1mm,before skip=6pt,after skip=6pt}
\newtcolorbox{code}[1][Code]{breakable,colback=black!4,colframe=black!60,
  fonttitle=\bfseries\small,title={#1},boxrule=0.5pt,arc=1mm,fontupper=\ttfamily\small,
  before skip=6pt,after skip=6pt}
% Blocs de code coloré (écrire le code sans accents : comme dans l'éditeur)
\newtcblisting{arduino}[1][Code Arduino]{listing only,breakable,colback=black!4,
  colframe=black!60,fonttitle=\bfseries\small,title={#1},boxrule=0.5pt,arc=1mm,
  before skip=6pt,after skip=6pt,
  listing options={language=C++,basicstyle=\ttfamily\small,columns=fullflexible,
  keepspaces=true,showstringspaces=false,commentstyle=\color{green!40!black}\itshape,
  keywordstyle=\color{blue!70!black}\bfseries}}
\newtcblisting{python}[1][Code Python]{listing only,breakable,colback=black!4,
  colframe=black!60,fonttitle=\bfseries\small,title={#1},boxrule=0.5pt,arc=1mm,
  before skip=6pt,after skip=6pt,
  listing options={language=Python,basicstyle=\ttfamily\small,columns=fullflexible,
  keepspaces=true,showstringspaces=false,commentstyle=\color{green!40!black}\itshape,
  keywordstyle=\color{blue!70!black}\bfseries}}
\newcommand{\entete}[2]{%
  \begin{tcolorbox}[colback=#1!10,colframe=#1!70!black,boxrule=0.8pt,arc=2mm]
  {\Large\bfseries #2}\hfill Nom : \rule{4cm}{0.3pt}\quad Date : \rule{2.2cm}{0.3pt}
  \end{tcolorbox}}

% Bandeau de titre : #1 titre, #2 semaine
\newcommand{\titreblock}[2]{%
\begin{tcolorbox}[colback=teal!10,colframe=teal!70!black,arc=2mm,boxrule=1pt]
\centering
{\LARGE\bfseries #1}\\[2mm]
{\large Robotique -- \gradelabel{} -- Plan de cours de 2 heures}\\[1mm]
Enseignant : Francis Ki \quad|\quad #2
\end{tcolorbox}}

% Tableau du minutage : 6 contenus (accroche, mini-leçon, guidée, autonome, mise en commun, exit ticket)
\newcommand{\minutage}[6]{%
\begin{tabularx}{\linewidth}{>{\bfseries}p{2.2cm}p{2.6cm}Y p{3cm}}
\toprule
Minutes & Étape & Contenu & Organisation\\
\midrule
0 -- 10 & Accroche & #1 & Classe entière\\
10 -- 30 & Mini-leçon & #2 & Classe entière\\
30 -- 60 & Pratique guidée & #3 & Collectif puis binômes\\
60 -- 95 & Pratique autonome & #4 & Individuel sur poste\\
95 -- 110 & Mise en commun & #5 & Classe entière\\
110 -- 120 & Exit ticket & #6 & Individuel sur poste\\
\bottomrule
\end{tabularx}}

% Étayage progressif : #1 aide 1, #2 aide 2, #3 aide 3
\newcommand{\aides}[3]{%
\textbf{Étayage progressif de l'aide.}
\begin{center}
\begin{tabularx}{\linewidth}{>{\bfseries}p{2.4cm}Y p{4.4cm}}
\toprule
Niveau d'aide & Contenu & Qui intervient\\
\midrule
Aide 1 & #1 & Le binôme seul\\
Aide 2 & #2 & Le binôme avec la carte d'aide\\
Aide 3 & #3 & L'enseignant\\
\bottomrule
\end{tabularx}
\end{center}}

% Les trois groupes de la pratique autonome : #1 tâche A, #2 tâche B, #3 tâche C
\newcommand{\groupes}[3]{%
\textbf{Consigne commune\,:} chaque élève ouvre le fichier de son groupe, le renomme avec son prénom et travaille seul sur son poste. L'enseignant passe d'abord auprès du groupe A pendant 10 min.

\medskip
\begin{tabularx}{\linewidth}{>{\bfseries}p{3cm}Y p{2.2cm}}
\toprule
Groupe & Tâche sur poste & Annexe\\
\midrule
A : en dessous du niveau & #1 & Annexe C\\
B : au niveau & #2 & Annexe D\\
C : au-dessus du niveau & #3 & Annexe E\\
\bottomrule
\end{tabularx}}

% Mise en commun : #1 groupe A, #2 groupe B, #3 groupe C, #4 conclusion
\newcommand{\miseencommun}[4]{%
\begin{etape}{teal}{5. Mise en commun}{95 -- 110 min}
\begin{enumerate}
 \item (2 min) Retour au calme, écrans mis en veille.
 \item (9 min) Trois passages courts de 3 min, avec valorisation :
 \begin{itemize}
  \item \textbf{Groupe A} : #1
  \item \textbf{Groupe B} : #2
  \item \textbf{Groupe C} : #3
 \end{itemize}
 \item (4 min) Conclusion : #4
\end{enumerate}
\end{etape}}

% Clôture : #1 question 1, #2 question 2
\newcommand{\cloture}[2]{%
\begin{etape}{gray}{6. Clôture : exit ticket}{110 -- 120 min}
Formulaire en ligne ou papier (annexe G) :
\begin{enumerate}
 \item #1
 \item #2
 \item Auto-évaluation : \emph{Je maîtrise / Je progresse / J'ai besoin d'aide.}
\end{enumerate}
Lire les réponses avant la séance suivante pour ajuster les groupes de besoin.
\end{etape}}

% Grille d'observation : #1 #2 #3 intitulés de colonnes
\newcommand{\grille}[3]{%
\subsection*{Grille d'observation (enseignant)}
\begin{tabularx}{\linewidth}{Y p{2.4cm} p{2.4cm} p{2.4cm} Y}
\toprule
Élève & #1 & #2 & #3 & Groupe suivant\\
\midrule
 & & & & \\ & & & & \\ & & & & \\ & & & & \\ & & & & \\
\bottomrule
\end{tabularx}}

% Signature de l'exit ticket
\newcommand{\autoeval}{Je me sens : \cb Je maîtrise \quad \cb Je progresse \quad \cb J'ai besoin d'aide}
.
\usepackage[T1]{fontenc}
\usepackage[utf8]{inputenc}
\usepackage{lmodern}
\usepackage[french]{babel}
\usepackage[margin=2cm,headheight=14pt]{geometry}
\usepackage[most]{tcolorbox}
\usepackage{xcolor,tabularx,booktabs,array,enumitem,fancyhdr,amssymb,tikz}
\usetikzlibrary{arrows.meta,positioning}

\pagestyle{fancy}
\fancyhf{}
\lhead{\small Robotique -- \gradelabel{} -- Novembre 2026}
\rhead{\small \leconhdr}
\cfoot{\thepage}
\setlist{itemsep=1pt,topsep=2pt,leftmargin=*}
\renewcommand{\arraystretch}{1.35}
\newcolumntype{Y}{>{\raggedright\arraybackslash}X}
\newcommand{\cb}{$\square$\,}
\newcommand{\fl}{$\rightarrow$}
\newcommand{\lig}{\noindent\rule{\linewidth}{0.3pt}\\[5mm]}

\newtcolorbox{etape}[3]{enhanced,breakable,colback=#1!6,colframe=#1!70!black,
  coltitle=white,fonttitle=\bfseries,title={#2\hfill #3},boxrule=0.8pt,arc=2mm,
  before skip=8pt,after skip=8pt}
\newtcolorbox{astuce}[1][Différenciation et besoins particuliers]{breakable,
  colback=yellow!12,colframe=yellow!50!black,fonttitle=\bfseries\small,
  title={#1},boxrule=0.5pt,arc=1mm,before skip=6pt,after skip=6pt}
\newtcolorbox{code}[1][Code]{breakable,colback=black!4,colframe=black!60,
  fonttitle=\bfseries\small,title={#1},boxrule=0.5pt,arc=1mm,fontupper=\ttfamily\small,
  before skip=6pt,after skip=6pt}
% Blocs de code coloré (écrire le code sans accents : comme dans l'éditeur)
\newtcblisting{arduino}[1][Code Arduino]{listing only,breakable,colback=black!4,
  colframe=black!60,fonttitle=\bfseries\small,title={#1},boxrule=0.5pt,arc=1mm,
  before skip=6pt,after skip=6pt,
  listing options={language=C++,basicstyle=\ttfamily\small,columns=fullflexible,
  keepspaces=true,showstringspaces=false,commentstyle=\color{green!40!black}\itshape,
  keywordstyle=\color{blue!70!black}\bfseries}}
\newtcblisting{python}[1][Code Python]{listing only,breakable,colback=black!4,
  colframe=black!60,fonttitle=\bfseries\small,title={#1},boxrule=0.5pt,arc=1mm,
  before skip=6pt,after skip=6pt,
  listing options={language=Python,basicstyle=\ttfamily\small,columns=fullflexible,
  keepspaces=true,showstringspaces=false,commentstyle=\color{green!40!black}\itshape,
  keywordstyle=\color{blue!70!black}\bfseries}}
\newcommand{\entete}[2]{%
  \begin{tcolorbox}[colback=#1!10,colframe=#1!70!black,boxrule=0.8pt,arc=2mm]
  {\Large\bfseries #2}\hfill Nom : \rule{4cm}{0.3pt}\quad Date : \rule{2.2cm}{0.3pt}
  \end{tcolorbox}}

% Bandeau de titre : #1 titre, #2 semaine
\newcommand{\titreblock}[2]{%
\begin{tcolorbox}[colback=teal!10,colframe=teal!70!black,arc=2mm,boxrule=1pt]
\centering
{\LARGE\bfseries #1}\\[2mm]
{\large Robotique -- \gradelabel{} -- Plan de cours de 2 heures}\\[1mm]
Enseignant : Francis Ki \quad|\quad #2
\end{tcolorbox}}

% Tableau du minutage : 6 contenus (accroche, mini-leçon, guidée, autonome, mise en commun, exit ticket)
\newcommand{\minutage}[6]{%
\begin{tabularx}{\linewidth}{>{\bfseries}p{2.2cm}p{2.6cm}Y p{3cm}}
\toprule
Minutes & Étape & Contenu & Organisation\\
\midrule
0 -- 10 & Accroche & #1 & Classe entière\\
10 -- 30 & Mini-leçon & #2 & Classe entière\\
30 -- 60 & Pratique guidée & #3 & Collectif puis binômes\\
60 -- 95 & Pratique autonome & #4 & Individuel sur poste\\
95 -- 110 & Mise en commun & #5 & Classe entière\\
110 -- 120 & Exit ticket & #6 & Individuel sur poste\\
\bottomrule
\end{tabularx}}

% Étayage progressif : #1 aide 1, #2 aide 2, #3 aide 3
\newcommand{\aides}[3]{%
\textbf{Étayage progressif de l'aide.}
\begin{center}
\begin{tabularx}{\linewidth}{>{\bfseries}p{2.4cm}Y p{4.4cm}}
\toprule
Niveau d'aide & Contenu & Qui intervient\\
\midrule
Aide 1 & #1 & Le binôme seul\\
Aide 2 & #2 & Le binôme avec la carte d'aide\\
Aide 3 & #3 & L'enseignant\\
\bottomrule
\end{tabularx}
\end{center}}

% Les trois groupes de la pratique autonome : #1 tâche A, #2 tâche B, #3 tâche C
\newcommand{\groupes}[3]{%
\textbf{Consigne commune\,:} chaque élève ouvre le fichier de son groupe, le renomme avec son prénom et travaille seul sur son poste. L'enseignant passe d'abord auprès du groupe A pendant 10 min.

\medskip
\begin{tabularx}{\linewidth}{>{\bfseries}p{3cm}Y p{2.2cm}}
\toprule
Groupe & Tâche sur poste & Annexe\\
\midrule
A : en dessous du niveau & #1 & Annexe C\\
B : au niveau & #2 & Annexe D\\
C : au-dessus du niveau & #3 & Annexe E\\
\bottomrule
\end{tabularx}}

% Mise en commun : #1 groupe A, #2 groupe B, #3 groupe C, #4 conclusion
\newcommand{\miseencommun}[4]{%
\begin{etape}{teal}{5. Mise en commun}{95 -- 110 min}
\begin{enumerate}
 \item (2 min) Retour au calme, écrans mis en veille.
 \item (9 min) Trois passages courts de 3 min, avec valorisation :
 \begin{itemize}
  \item \textbf{Groupe A} : #1
  \item \textbf{Groupe B} : #2
  \item \textbf{Groupe C} : #3
 \end{itemize}
 \item (4 min) Conclusion : #4
\end{enumerate}
\end{etape}}

% Clôture : #1 question 1, #2 question 2
\newcommand{\cloture}[2]{%
\begin{etape}{gray}{6. Clôture : exit ticket}{110 -- 120 min}
Formulaire en ligne ou papier (annexe G) :
\begin{enumerate}
 \item #1
 \item #2
 \item Auto-évaluation : \emph{Je maîtrise / Je progresse / J'ai besoin d'aide.}
\end{enumerate}
Lire les réponses avant la séance suivante pour ajuster les groupes de besoin.
\end{etape}}

% Grille d'observation : #1 #2 #3 intitulés de colonnes
\newcommand{\grille}[3]{%
\subsection*{Grille d'observation (enseignant)}
\begin{tabularx}{\linewidth}{Y p{2.4cm} p{2.4cm} p{2.4cm} Y}
\toprule
Élève & #1 & #2 & #3 & Groupe suivant\\
\midrule
 & & & & \\ & & & & \\ & & & & \\ & & & & \\ & & & & \\
\bottomrule
\end{tabularx}}

% Signature de l'exit ticket
\newcommand{\autoeval}{Je me sens : \cb Je maîtrise \quad \cb Je progresse \quad \cb J'ai besoin d'aide}


\begin{document}
\titreblock{Feu tricolore automatisé}{Semaine du 16 au 20 novembre 2026}

%==================================================================
\section{Informations générales}
\begin{tabularx}{\linewidth}{>{\bfseries}p{3.6cm}Y}
\toprule
Titre & Feu tricolore automatisé\\
Classe / effectif & Grade 9, 11 à 20 élèves, niveaux hétérogènes, élèves à besoins particuliers\\
Durée & 2 heures (120 min)\\
Lieu & Salle informatique, un poste par élève\\
Matériel et logiciels &
\begin{itemize}
 \item Un poste par élève avec navigateur web : \textbf{Tinkercad Circuits} (simulation).
 \item Arduino IDE sur les postes pour vérifier le code.
 \item Facultatif : une carte Arduino, trois LEDs (rouge, orange, verte) et trois résistances de 220~$\Omega$ par binôme.
 \item Un outil de dessin d'algorigramme (diapositives, ou papier et crayon).
\end{itemize}\\
Pré-requis &
Séances précédentes : \texttt{pinMode}, \texttt{digitalWrite}, \texttt{delay}, chenillard de trois LEDs.\\
\bottomrule
\end{tabularx}

\begin{astuce}[Préparation avant le cours (environ 40 min)]
\begin{enumerate}
 \item Préparer un circuit modèle partagé avec trois LEDs sur les broches 8 (rouge), 9 (orange), 10 (verte).
 \item Préparer un second circuit à six LEDs (deux feux) pour les groupes B et C.
 \item Préparer les fichiers des annexes C à E et le formulaire de sortie.
\end{enumerate}
\end{astuce}

%==================================================================
\section{Objectifs d'apprentissage}
\begin{tcolorbox}[colback=gray!8,colframe=gray!60!black,arc=1mm,boxrule=0.5pt]
\textbf{Objectifs du programme}
\begin{itemize}
 \item Concevoir un système automatisé simple.
 \item Utiliser des séquences de temporisation.
\end{itemize}
\textbf{À la fin de la séance, l'élève est capable\,:}
\begin{itemize}
 \item de décrire le cycle d'un feu tricolore (états et durées) dans un tableau ou un algorigramme\,;
 \item de traduire ce cycle en programme Arduino avec des temporisations\,;
 \item de tester le programme en simulation et de corriger une durée ou un ordre.
\end{itemize}
\textbf{Compétences numériques\,:} représenter un cycle, tester en simulation, comparer le comportement attendu et observé.
\end{tcolorbox}

%==================================================================
\section{Déroulé minuté (120 min)}
\minutage{Observation d'un carrefour et de son feu}{Cycle, états et algorigramme}{Un feu à trois LEDs en collectif, puis un feu piéton en binômes}{Trois niveaux de besoin}{Valoriser chaque groupe}{Trois questions sur poste}

%------------------------------------------------------------------
\begin{etape}{orange}{1. Accroche : le carrefour}{0 -- 10 min}
\begin{enumerate}
 \item (5 min) Projeter la photo d'un carrefour à deux feux. \emph{Quelles couleurs et dans quel ordre ? Pourquoi un moment où tous les feux sont rouges ?}
 \item (5 min) Les élèves miment le cycle d'un feu avec trois cartes colorées. Question-problème : \textbf{«~Comment faire fonctionner un feu tout seul, sans personne pour changer les couleurs ?~»}
\end{enumerate}
\end{etape}

%------------------------------------------------------------------
\begin{etape}{blue}{2. Mini-leçon : un cycle d'états}{10 -- 30 min}
\textbf{Support\,:} annexe A. Quatre temps de 5 min.

\textbf{Temps 1 : le besoin et le cycle.} Un système automatisé répète un cycle. Le feu passe par trois états :
\begin{center}
\begin{tikzpicture}[box/.style={draw,rounded corners,thick,minimum height=1cm,text width=2.4cm,align=center,font=\small},>={Stealth[length=3mm]},node distance=0.9cm]
\node[box,fill=green!25] (v) {Vert\\ 5 s};
\node[box,fill=orange!30,right=of v] (o) {Orange\\ 2 s};
\node[box,fill=red!25,right=of o] (r) {Rouge\\ 5 s};
\draw[->,thick] (v)--(o); \draw[->,thick] (o)--(r);
\draw[->,thick] (r) to[bend left=35] (v);
\end{tikzpicture}
\end{center}

\textbf{Temps 2 : le tableau des états.}
\begin{center}
\begin{tabularx}{\linewidth}{>{\bfseries}p{2cm}c c c Y}
\toprule
État & Rouge (8) & Orange (9) & Vert (10) & Durée\\
\midrule
1 & éteinte & éteinte & allumée & 5000 ms\\
2 & éteinte & allumée & éteinte & 2000 ms\\
3 & allumée & éteinte & éteinte & 5000 ms\\
\bottomrule
\end{tabularx}
\end{center}

\textbf{Temps 3 : le programme.}
\begin{arduino}[Feu tricolore]
void setup() {
  pinMode(8, OUTPUT);   // rouge
  pinMode(9, OUTPUT);   // orange
  pinMode(10, OUTPUT);  // vert
}
void loop() {
  digitalWrite(8, LOW);  digitalWrite(9, LOW);  digitalWrite(10, HIGH);
  delay(5000);
  digitalWrite(10, LOW); digitalWrite(9, HIGH);
  delay(2000);
  digitalWrite(9, LOW);  digitalWrite(8, HIGH);
  delay(5000);
}
\end{arduino}

\textbf{Temps 4 : de l'état au code.} Chaque ligne du tableau devient un bloc d'instructions suivi d'un \texttt{delay}. En classe on peut réduire les durées (3, 1, 3 s) pour tester plus vite.

\begin{astuce}
Utiliser les couleurs réelles dans le tableau. Pour les élèves qui ont du mal à lire le code, fournir le tableau des états avec les lignes du programme associées. Autoriser des durées réduites pour les essais.
\end{astuce}
\end{etape}

%------------------------------------------------------------------
\begin{etape}{green}{3. Pratique guidée}{30 -- 60 min}
\textbf{Fiche\,:} annexe B.

\textbf{3a. Un feu à trois LEDs en collectif (12 min).} Compléter ensemble le tableau des états, puis écrire le programme et le simuler.

\textbf{3b. Un feu piéton en binômes (18 min).} Ajouter un feu piéton (rouge et vert) synchronisé avec le feu des voitures : quand les voitures ont le vert, le piéton a le rouge, et inversement. Alternance des rôles.

\aides{Le tableau des états du 3a et le programme du feu}{Questions guides : \emph{Dans quel état le piéton peut-il traverser ? Quelles LEDs doivent être allumées ?}}{L'enseignant remplit une ligne du tableau avec le binôme, qui poursuit}

\begin{astuce}
Écrire en commentaire ce que fait chaque bloc d'instructions. Tester état par état avant d'assembler le cycle complet.
\end{astuce}
\end{etape}

%------------------------------------------------------------------
\begin{etape}{violet}{4. Pratique autonome différenciée}{60 -- 95 min}
\groupes{Programme du feu tricolore à compléter : instructions manquantes signalées par des tirets, tableau des états déjà rempli, banque de valeurs. Simuler et comparer avec le tableau.}{Programmer un carrefour à deux feux croisés (six LEDs) avec un état où tous les feux sont rouges pendant 1 s, après avoir rempli le tableau des états.}{Concevoir et programmer un carrefour avec feux voitures et feux piétons, en structurant le code avec des fonctions (\texttt{void feuAvert()}, \texttt{void feuRouge()}), et ajouter un mode nuit (orange clignotant).}

\medskip
\textbf{Plan B si la simulation ne fonctionne pas\,:} vérifier le tableau des états sur papier et le programme avec l'outil de vérification de l'IDE.
\end{etape}

%------------------------------------------------------------------
\miseencommun{un élève présente son feu et lit son tableau des états. Valoriser la correspondance entre le tableau et le programme.}{un binôme montre le carrefour à deux feux et explique l'état «~tout rouge~». Valoriser la sécurité.}{un élève présente son code structuré en fonctions. Valoriser la lisibilité.}{\textit{un système automatisé est un cycle d'états avec des durées ; on le décrit avant de le programmer.}}

\cloture{Dans quel ordre s'allument les couleurs d'un feu tricolore ?}{Pourquoi prévoit-on un court moment où tous les feux sont rouges dans un carrefour ?}

%==================================================================
% ANNEXES
%==================================================================
\newpage
\section*{Annexe A : Support résumé de la mini-leçon}
\entete{blue}{Feu tricolore automatisé}

\begin{tcolorbox}[colback=yellow!10,colframe=yellow!50!black]
Un système automatisé répète un \textbf{cycle d'états}. On décrit chaque état (LEDs allumées) et sa \textbf{durée}.
\end{tcolorbox}

\begin{tabularx}{\linewidth}{>{\bfseries}p{2cm}c c c Y}
\toprule
État & Rouge (8) & Orange (9) & Vert (10) & Durée\\
\midrule
1 & éteinte & éteinte & allumée & 5000 ms\\
2 & éteinte & allumée & éteinte & 2000 ms\\
3 & allumée & éteinte & éteinte & 5000 ms\\
\bottomrule
\end{tabularx}

\begin{arduino}[Structure du programme]
void loop() {
  // etat 1 : vert
  // etat 2 : orange
  // etat 3 : rouge
}
\end{arduino}

%==================================================================
\newpage
\section*{Annexe B : Fiche de pratique guidée}
\entete{green}{Pratique guidée : le feu tricolore}

\subsection*{3a. Tableau des états des voitures}
\begin{tabularx}{\linewidth}{>{\bfseries}p{2cm}c c c Y}
\toprule
État & Rouge (8) & Orange (9) & Vert (10) & Durée (ms)\\
\midrule
1 & & & & \\ 2 & & & & \\ 3 & & & & \\
\bottomrule
\end{tabularx}

\subsection*{3b. Feu piéton (rouge broche 11, vert broche 12)}
\begin{tabularx}{\linewidth}{>{\bfseries}p{2cm}Y Y Y}
\toprule
État voiture & Piéton rouge (11) & Piéton vert (12) & Pourquoi ?\\
\midrule
Vert & & & \\ Orange & & & \\ Rouge & & & \\
\bottomrule
\end{tabularx}

\begin{tcolorbox}[colback=yellow!10,colframe=yellow!50!black,title={Cartes d'aide à utiliser dans l'ordre},fonttitle=\bfseries\small]
\textbf{Aide 1 :} regarde le tableau des états du 3a.\\
\textbf{Aide 2 :} Dans quel état le piéton peut-il traverser ? Quelles LEDs doivent être allumées ?\\
\textbf{Aide 3 :} demande à l'enseignant de remplir une ligne avec toi.
\end{tcolorbox}

%==================================================================
\newpage
\section*{Annexe C : Niveau A (en dessous du niveau)}
\entete{orange}{Niveau A : je complète le feu}

\textbf{Banque :} HIGH, LOW, 5000, 2000.

\begin{arduino}[Programme à compléter]
void setup() {
  pinMode(8, OUTPUT);
  pinMode(9, OUTPUT);
  pinMode(10, OUTPUT);
}
void loop() {
  digitalWrite(8, LOW);  digitalWrite(9, LOW);  digitalWrite(10, _____);
  delay(_____);
  digitalWrite(10, LOW); digitalWrite(9, _____);
  delay(_____);
  digitalWrite(9, LOW);  digitalWrite(8, _____);
  delay(5000);
}
\end{arduino}

\subsection*{Activité : vérifie avec le tableau}
\begin{tabularx}{\linewidth}{>{\bfseries}p{2cm}Y Y}
\toprule
État & LED allumée & Durée\\
\midrule
1 & & \\ 2 & & \\ 3 & & \\
\bottomrule
\end{tabularx}

%==================================================================
\newpage
\section*{Annexe D : Niveau B (au niveau)}
\entete{blue}{Niveau B : carrefour à deux feux}

\textbf{Feu A :} rouge 8, orange 9, vert 10. \textbf{Feu B :} rouge 11, orange 12, vert 13.

\begin{tabularx}{\linewidth}{>{\bfseries}p{1.4cm}c c c c c c Y}
\toprule
État & A rouge & A orange & A vert & B rouge & B orange & B vert & Durée (ms)\\
\midrule
1 & & & & & & & \\ 2 & & & & & & & \\ 3 & & & & & & & \\ 4 & & & & & & & \\ 5 & & & & & & & \\ 6 & & & & & & & \\
\bottomrule
\end{tabularx}

\subsection*{Questions}
\begin{enumerate}
 \item Pourquoi prévoit-on un état où A et B sont rouges ensemble ? \dotfill
 \item Combien de secondes dure un cycle complet ? \dotfill
\end{enumerate}

%==================================================================
\newpage
\section*{Annexe E : Niveau C (au-dessus du niveau)}
\entete{violet}{Niveau C : carrefour structuré avec mode nuit}

\textbf{Mission :} programme un carrefour avec des fonctions.

\begin{arduino}[Aide : structure]
void feuVert()   { /* instructions */ }
void feuOrange() { /* instructions */ }
void feuRouge()  { /* instructions */ }

void loop() {
  feuVert();
  feuOrange();
  feuRouge();
}
\end{arduino}

\begin{enumerate}
 \item Écris le programme complet avec un mode nuit (orange clignotant, 500 ms) déclenché par une variable \texttt{nuit} : \lig \lig \lig
 \item Pourquoi les fonctions rendent-elles le programme plus facile à modifier ? \dotfill
\end{enumerate}

\textbf{Défi bonus :} ajoute un compteur de cycles (\texttt{int cycle}) et affiche-le avec \texttt{Serial.println} dans le moniteur série.

%==================================================================
\newpage
\section*{Annexe F : Corrections et attendus (enseignant)}

\subsection*{Annexe B}
États : vert (rouge 0, orange 0, vert 1, 5000 ms) ; orange (0, 1, 0, 2000) ; rouge (1, 0, 0, 5000). Piéton : voiture verte $\to$ piéton rouge ; voiture orange $\to$ piéton rouge ; voiture rouge $\to$ piéton vert.

\subsection*{Annexe C}
HIGH, 5000, HIGH, 2000, HIGH. États : vert, orange, rouge.

\subsection*{Annexe D}
Six états : A vert / B rouge ; A orange / B rouge ; tout rouge (1 s) ; A rouge / B vert ; A rouge / B orange ; tout rouge (1 s). L'état «~tout rouge~» laisse le temps aux véhicules de dégager le carrefour. Durée d'un cycle : somme des durées (par exemple 5 + 2 + 1 + 5 + 2 + 1 = 16 s).

\subsection*{Annexe E}
Mode nuit : \texttt{if (nuit) \{ digitalWrite(9, HIGH); delay(500); digitalWrite(9, LOW); delay(500); \} else \{ feuVert(); feuOrange(); feuRouge(); \}}. Les fonctions regroupent les instructions et évitent de recopier.\\
\textbf{Critères de réussite :} trois fonctions appelées dans le bon ordre, mode nuit fonctionnel.

%==================================================================
\newpage
\section*{Annexe G : Exit ticket et grille d'observation}
\entete{gray}{Exit ticket : feu tricolore}

\begin{enumerate}
 \item Dans quel ordre s'allument les couleurs d'un feu tricolore ?\\ \lig
 \item Pourquoi prévoit-on un court moment où tous les feux sont rouges dans un carrefour ?\\ \lig
 \item \autoeval
\end{enumerate}
\textbf{Attendus :} 1. Vert, orange, rouge. 2. Pour laisser les véhicules finir de traverser avant le vert de l'autre feu.

\bigskip
\grille{Tableau des états}{Programme fonctionnel}{Simulation réussie}

\end{document}
